\chapter*{General Introduction}
\addcontentsline{toc}{chapter}{General Introduction}

\noindent
Enterprise competitiveness increasingly depends on how well an organization manages the relationships around it. A company such as \reportcompany{} works with corporate customers, technology and service suppliers, marketing organizations, and engineering schools. Each of these relationships carries its own contacts, commercial terms, shared projects, documents, and obligations. When this network grows, managing it through email threads, spreadsheets, and shared drives becomes error-prone. Renewals slip, warning signs stay hidden until attrition is already under way, and teams lose the common view they need to decide where to focus.

\projectname{} was designed and built during a master's final-year internship to replace that fragmented tooling with one governed platform. It centralizes the full partnership lifecycle and the project portfolio behind a secure, role-aware interface, and it adds an intelligence layer that turns operational data into grounded decision support.

\medskip

The platform is organized around three surfaces. The public website presents the company's partnership offer and lets an organization submit a partnership request that enters the internal review queue directly. The employee portal is the operational hub for Capgemini staff, giving access to partners, communications, documents, project delivery, business-intelligence dashboards, and the AI assistant. The partner portal gives each external organization a secure self-service space adapted to its category, covering customers, suppliers, marketing partners, and universities.

The intelligence layer is the distinguishing contribution of the project. \projectname{} implements an agentic assistant built on a streaming conversational interface and a typed, validated tool set, where every answer is grounded in real data from the live database. It combines deterministic partner and project scoring with semantic document retrieval, and it produces executive reports that analysts can export for real use.

\medskip

These choices are driven by a single question that runs through the report:

\begin{quote}
\itshape
How can a unified, role-aware platform cover the full operational lifecycle of enterprise partnerships and projects while embedding a governed agentic AI layer whose outputs are grounded, measurable, and trustworthy in a production context?
\end{quote}

\medskip

\noindent
The remainder of this report is organized into seven chapters.

\begin{itemize}
  \item \textbf{Chapter 1, General Presentation and Methodology}, introduces the host company, the business context of enterprise partnership management, the project objectives, the positioning of \projectname{} against existing solutions, and the agile method followed during the internship.

  \item \textbf{Chapter 2, Planning and Architecture}, defines the functional and non-functional requirements, the product backlog and sprint plan, the use-case and data models, the overall application architecture, and the technical environment.

  \item \textbf{Chapter 3, Access, Identity and Auditability}, covers authentication, the role-based access-control model, session handling, and the audit trail that records every sensitive operation across the platform.

  \item \textbf{Chapter 4, Partnership Lifecycle Management and Customers}, presents the partner registry and status workflow, the contact, offer, and event modules, secure document handling, notifications, and the public request intake and review pipeline.

  \item \textbf{Chapter 5, Project Delivery, Operations and Business Intelligence}, covers project delivery, the human-resources workflows, and the data-warehouse-backed analytics dashboard.

  \item \textbf{Chapter 6, Agentic AI Harness and Generative Interface}, examines the conversational agent in depth: its validated tool set, the streaming protocol, multi-step plan execution and the live plan interface, in-line charts and tables, and the report-generation pipeline.

  \item \textbf{Chapter 7, Knowledge Retrieval, Scoring and Observability}, covers document retrieval, the partner and project scoring models, churn-risk prediction, recommendations, and the evaluation and observability framework used to keep AI output measurable.
\end{itemize}

\medskip

\noindent
A \textbf{General Conclusion} closes the report with a synthesis of the achieved results, an assessment against the initial objectives, and perspectives for future development and industrialization.
